% Part: incompleteness
% Chapter: arithmetization-syntax
% Section: substitution

\documentclass[../../../include/open-logic-section]{subfiles}

\begin{document}

\olfileid{inc}{art}{sub}
\olsection{आदेशन}

आदेशन म्हणजे !!a{formula}~$!A$ मध्ये !!a{variable}~$u$ च्या
सर्व मुक्त आस्थितींच्या जागी पद~$t$ ठेवण्याची क्रिया आहे;
ती $\Subst{!A}{t}{u}$ अशी लिहितात, हे आठवा. !!{variable}s,
!!{formula}s आणि पदे यांच्या गोडेल संख्यांवर ही क्रिया
केली, तर ती आदिम पुनरावर्ती असते.

\begin{prop}\ollabel{prop:subst-primrec}
पुढील गुणधर्म असलेले आदिम पुनरावर्ती फलन $\fn{Subst}(x, y, z)$
अस्तित्वात आहे:
\[
\fn{Subst}(\Gn{!A}, \Gn{t}, \Gn{u}) = \Gn{\Subst{!A}{t}{u}}.
\]
\end{prop}

\begin{proof}
आता $\fn{hSubst}$ हे फलन पुढीलप्रमाणे आदिम पुनरावर्तनाने
परिभाषित करता येते:
\begin{multline*}
\begin{aligned}
\fn{hSubst}(x, y, z, 0) & = \emptyseq \\
\fn{hSubst}(x, y, z, i+1) & =
\end{aligned}\\
\begin{cases}
\fn{hSubst}(x, y, z, i) \concat y & \text{जर $\fn{FreeOcc}(x, z, i)$ खरे असेल तर} \\
\fn{append}(\fn{hSubst}(x, y, z, i), (x)_{i}) & \text{अन्यथा.}
\end{cases}
\end{multline*}
आता $\fn{Subst}(x, y, z)$ ची व्याख्या
$\fn{hSubst}(x, y, z, \len{x})$ अशी करता येते.
\end{proof}

\begin{prop}
\ollabel{prop:free-for}
गोडेल संख्या~$y$ असलेले पद, गोडेल संख्या~$x$ असलेल्या सूत्रात
गोडेल संख्या~$z$ असलेल्या चरासाठी !!{free for} असेल,
तर आणि तरच खरे असणारा $\fn{FreeFor}(x, y, z)$ हा संबंध
आदिम पुनरावर्ती आहे.
\end{prop}

\begin{proof} सराव. \end{proof}

\begin{prob}
\olref[inc][art][sub]{prop:free-for} सिद्ध करा.
\end{prob}

\end{document}
